\input{estilo}
% La pizarra: tres particiones (1, 2, 3) como tiras largas; C₁ lee la 1 y la
% 2, C₂ (tachado) la 3; C₁ habla con el coordinador (el ida y vuelta marcado
% "1") y escribe en ZooKeeper.
\begin{document}
\begin{figura}\begin{tikzpicture}[q/.style={caja, minimum width=11mm, minimum height=11mm}]
  \foreach \i/\y in {1/2.0, 2/0.9, 3/-0.9} {
    \node at (-0.6,\y) {\i};
    \draw (0,\y-0.3) rectangle (5.4,\y+0.3);
  }
  \node[q] (c1) at (10.2,1.5) {\var{C}\textsubscript{1}};
  \node[q] (c2) at (10.2,-0.9) {\var{C}\textsubscript{2}}; \tachar{c2}
  \draw[msg] (c1.west) to[out=190, in=-10] (5.55,2.1);
  \draw[msg] (5.55,0.85) to[out=0, in=200] (c1.south west);
  \draw[msg] (5.55,-1.0) to[out=0, in=180] (c2.west);
  \node[caja, minimum width=30mm, minimum height=10mm] (k) at (7.6,3.6) {coordinador};
  \draw[msg] (c1.north west) to[bend left=30] (k.south east);
  \draw[msg] (k.south east) ++(-0.3,0) to[bend right=40] (c1.north west);
  \node at (9.75,2.9) {1};
  \pic at (12.6,3.4) {base};
  \node[anchor=north] at (12.6,3.3) {ZooKeeper};
  \draw[msg] (c1.north east) -- (12.0,3.9);
\end{tikzpicture}\end{figura}
\end{document}
